% thesis.tex
%
% This file is root file for an example thesis written using the
% IIT Bombay LaTeX Style file.
% Created by Amey Karkare (21 June 2007)
%
% It is provided without warranty on an AS IS basis.

%=====================================================================
% Read: http://www.cse.iitb.ac.in/karkare/iitbthesis/
%    FAQ.txt     for frequently asked quetions
%    Changes.txt for changes
%    README      for more information
%=====================================================================

%=====================================================================
% DOCUMENT STYLE
%=====================================================================
% IITB PhD Thesis format default settings are:
%   12pt, one-sided printing on a4 size paper
\documentclass[openright,oneside]{iitkthesis}
% For two-sided printing, with Chapter starting on odd-numbered pages,
% use the following line instead:  
%\documentclass[openright,twoside]{iitkthesis}

%=====================================================================
% OPTIONAL PACKAGES
%=====================================================================
% To include optional packages, use the \usepackage command.
% For e.g., The package epsfig is used to bring in the Encapsulated
%    PostScript figures into the document.
%    The package times is used to change the fonts to Times Roman;
%=====================================================================

\usepackage{amsmath}
\usepackage{amsthm}
\usepackage{amssymb}
\usepackage{lipsum}
\usepackage[hidelinks]{hyperref}
\usepackage{microtype}
\usepackage[font=small,labelfont=bf]{caption}
\usepackage{fancyref}
\usepackage{color,soul} % for highlights
\usepackage{enumerate}
\usepackage{tabularx}
\usepackage{graphicx}
\graphicspath{{figures/}}

%to add signed cerificate
\usepackage{pdfpages}

% to write the date in the desired format
\usepackage{datetime}

%=====================================================================
%  Single counter for theorems and theorem-like environments:
%=====================================================================
\theoremstyle{plain}
\newtheorem{theorem}{Theorem}
\newtheorem{lemma}[theorem]{Lemma}
\newtheorem{corollary}[theorem]{Corollary}
\theoremstyle{definition}
\newtheorem{definition}[theorem]{Definition}
\newtheorem{example}[theorem]{Example}
\theoremstyle{remark}
\newtheorem*{remark}{Remark}
\newtheorem{claim}[theorem]{Claim}


%Other macros


\setlength{\headheight}{15pt}

%\expandafter\def\csname PY@tok@err\endcsname{}

%=====================================================================
% End of Preamble, start of document
%

\begin{document}

%=====================================================================
% Include the prelude for Title page, abstract, table of contents, etc
% You need to modify it to contain your details
% prelude.tex
%   - titlepage
%   - dedication (optional)
%   - approval sheet
%   - course certificate
%   - table of contents, list of tables and list of figures
%   - nomenclature
%   - abstract
%============================================================================


\clearpage\pagenumbering{roman}  % This makes the page numbers Roman (i, ii, etc)



% TITLE PAGE
%   - define \title{} \author{} \date{}
\title{Title}
\author{Name}
\newdateformat{monthyear}{\monthname[\THEMONTH], \THEYEAR}
\date{\monthyear\today}

%  - Roll number, required for title page, approval sheet, and
%    certificate of course work 
\rollnum{RollNo} 

%   - The default degree is ``Doctor of Philosophy''
%     (unless the document style msthesis is specified
%      and then the default degree is ``Master of Science'')
%     Degree can be changed using the command \iitbdegree{}
\iitbdegree{Master of Technology}

%   - The default report type is preliminary report.
%      * for a PhD thesis, specify \thesis
\thesis
%      * for a M.Tech./M.Phil./M.Des./M.S. dissertation, specify \dissertation
%\dissertation
%      * for a DIIT/B.Tech./M.Sc.project report, specify \project
%\project
%      * for any other type, use  \reporttype{}
%\reporttype{ReportType}

%   - The default department is ``Unknown Department''
%     The department can be changed using the command \department{}
\department{Computer Science \& Engineering}

%    - Set the guide's name
\setguide{Guide}
\setguidedept{Department of Computer Science \& Engineering}

%   - once the above are defined, use \maketitle to generate the titlepage
\maketitle

%--------------------------------------------------------------------%
% CERTIFICATE
%     The first page after the title page.
\makecertificate
%\includepdf[pages=-]{thesis_front.pdf}

%--------------------------------------------------------------------%
% COPYRIGHT PAGE
%   - To include a copyright page use \copyrightpage
% \copyrightpage

%--------------------------------------------------------------------%
% ABSTRACT
\cleardoublepage
\phantomsection \label{abstract}
\addcontentsline{toc}{chapter}{Abstract}
\begin{abstract}
Abcd abstract

\end{abstract}

%--------------------------------------------------------------------%
% DEDICATION
%   Dedications, if any.
\begin{dedication}
To my family
\end{dedication}

% Acknowledgements
\begin{acknowledgments}

Abcd Acknol

\end{acknowledgments}

%--------------------------------------------------------------------%
% CONTENTS, TABLES, FIGURES
\setcounter{secnumdepth}{3}
\phantomsection \label{toc}
\addcontentsline{toc}{chapter}{Contents}
\tableofcontents
%\listoftables

\cleardoublepage
\phantomsection \label{listoffig}
\addcontentsline{toc}{chapter}{List of Figures}
\listoffigures

\cleardoublepage\pagenumbering{arabic} % Make the page numbers Arabic (1, 2, etc)


%=====================================================================
% Include the technical part of the report
%% \include{chap_intro}             % Chapter 1: Introduction
%% \include{chap_others}            % Other chapters as required
%% \include{chap_conclusions}       % Finally the summary & conclusions

%=====================================================================
% APPENDIX
%  Appendices, if any, must precede the cited literatures.
%  Appendices shall be numbered in Roman Capitals (e.g. Appendix IV)

%% \appendix
%% \include{appendix_something}          

%=====================================================================
% PUBLICATIONS
%  publications if any may be listed after the literature cited.
%% \include{mypubs}

%=====================================================================
% ACKNOWLEDGMENTS
%   This is the last item in the thesis. It should be signed by
%   author, with date.


\chapter{Introduction}
\label{chap:intro}

First Cahpter.

\section{Graph Reachability}
First Section


%\input{definitions}
%\input{preworks}
%\input{result}
%\input{conclusions}

\bibliographystyle{plainurl}% the recommended bibstyle
\cleardoublepage
\phantomsection \label{referen}
\addcontentsline{toc}{chapter}{Bibliography}
%\renewcommand\bibname{References}
\bibliography{citations}

\end{document}
